% ============================================================
%           CHAPTER 9: CONCLUSION & FUTURE WORK
% ============================================================
\chapter{Conclusion \& Future Work}

\section{Conclusion}

This exploratory project successfully conceptualized, constructed, and profiled a full-scale scenario testing framework heavily leveraging the co-simulation paradigm between SUMO and Unity3D. The systematic integration established herein unequivocally proves that high-fidelity 3D graphics can be seamlessly synchronized with rigorous, mathematically validated background traffic flow logic.

The core technical milestones achieved include:
\begin{enumerate}[label={\colorbox{PrimaryNavy}{\makebox[1.2em][c]{\color{white}\bfseries\arabic*}}}, leftmargin=3em, itemsep=8pt]
    \item \textbf{High-Frequency Bidirectional Sync:} Implementation of a NetMQ ZeroMQ asynchronous architecture bypassing legacy TraCI TCP bounds. This infrastructure delivered consistent bidirectional state transfers under $5\text{ms}$.
    \item \textbf{Algorithmic Terrain Parsing:} Execution of C\# routines that ingest raw coordinate dictionaries exported by SUMO to procedurally extrude realistic URP meshes, enabling massive reduction in manual aesthetic asset creation.
    \item \textbf{Ego-Physics Preservation:} Creation of robust interaction models separating rigid-body user interaction from static interpolation algorithms, achieving $112$ FPS geometric mean thresholds mandatory for real-time visual stability.
    \item \textbf{Hardware Integration Structure:} Successful bridging of user input peripherals directly into the Unity Execution cycle, thereby closing the loop of human-in-the-loop engineering.
\end{enumerate}

Ultimately, the driving simulator constructed offers a cost-effective, hyper-scalable, and deeply realistic tool capable of immediately collecting Time-to-Collision analytics on human factors research prior to engaging physical vehicles on physical asphalt.

\section{Future Work}

As the domain of traffic behavior simulation expands toward accommodating Artificial Intelligence endpoints, the primary framework generated by this project must scale accordingly. The recommended future engineering endpoints include:

\begin{itemize}[label={\color{AccentTeal}$\blacktriangleright$}, leftmargin=2em, itemsep=6pt]
    \item \textbf{Protocol Buffers Migration:} Exterminating JSON payload structures inherently caps agent scaling. Developing a lightweight protobuf binary wrapper across the NetMQ boundary will extend hardware limitations from $150$ vehicles to $>10,000$ simulated actors per scenario block.
    \item \textbf{Integration of Advanced Sensor Systems:} While human drivers react organically, validating Autonomous Vehicles (AV) algorithms necessitates replicating LiDAR clouds, stereo cameras, and Ultrasonic ranges. Building shader-level raycast buffers inside Unity to spoof data for ROS/Autoware stacks represents the apex application of this simulator.
    \item \textbf{Behavioral Data Science and Cloud Parsing:} Given that trajectory maps output continuously to CSV strings, connecting the telemetry port to a cloud-based Python ingestion server (e.g., Pandas processing) would permit fully automated surrogate safety measure (PET, TTC, DRAC) graph generation milliseconds after the simulation stops running.
    \item \textbf{Multi-Ego Instantiation:} Architecting network synchronization to permit multiple physical participants—for example, one user commanding an Ego Vehicle via desktop and another distinct VR subject commanding an Ego Pedestrian crossing the street simultaneously—would enable revolutionary interaction analyses entirely disconnected from generic AI behavioral algorithms.
\end{itemize}
